\chapter{Conclusion}\label{ch:conclusion}
\section{Summary of Findings}
The project delivers a local Vietnamese legal RAG prototype and a separate benchmark workflow. It combines hybrid retrieval, optional graph expansion and reranking, hierarchical context, pre-generation sufficiency evaluation, source-derived citations and post-generation verification. Models remain pretrained and the graph uses source-backed document references.

The development study covers 45 questions, three topics and five configurations, with 225 complete inference records. Each method completes 43 questions and returns insufficient evidence for two; no inference failure or timeout is recorded. Retrieval Recall@5 ranges from 0.9667 to 0.9778 and nDCG@5 from 0.9508 to 0.9614. Precision@5 ranges from 0.2133 to 0.2178 against a label-defined ceiling of 0.2267. These results support strong retrieval of the specified labels on this development set, not universal legal-question accuracy.

The optional components do not establish a statistically significant quality advantage under the reported Holm-adjusted comparisons. B0 has the lowest observed all-outcome mean inference latency at 21.76 seconds; B4 has the highest at 24.28 seconds. B4's additional rounds recover no extra completed answers, and its final answers match B3 for all 45 questions. No explicit graph-relation block survives into packed context in the recorded pilot, so a graph-grounded multi-hop benefit is not demonstrated.

\section{Contributions and Reflections on the Project}
The engineering contribution is a pipeline with visible evidence boundaries and failure states. Controlled method switches, source-resolved qrels, separate scoring artifacts, configuration identities, telemetry and paired statistics make its results inspectable, including when additions do not outperform the baseline.

The methodological lesson is that an operational benchmark is not automatically a valid correctness measure. This run is complete and arithmetically auditable, yet its judge has truncation errors and incorrect entailment decisions. Low faithfulness means cannot be converted directly into a hallucination rate.

\section{Limitations and Future Work}
The current set has informed development, contains only answerable questions, and awaits independent legal review. Its evidence labels are narrow, its task mix is dominated by lookup and multi-clause questions, and it has only one repeat per configuration. Hardware and LLM quantization records are incomplete. Neither a current-validity service nor a production web application is established by this implementation.

Future work should calibrate the judge against reviewed labels and rescore the frozen answers. A larger held-out test should include graph-dependent, multi-document, ambiguous and unanswerable requests. Repeated runs with complete environment records are needed before stronger comparative or deployment claims.
